%%%%%%%%%%%%%%%%%%%%%%%%%%%%%%%%%%%%%%%%%%%%%%%%%%%%%%%%%%%%%%
%%%%%%%%%%%%%%%%%% MA-0609 Teoría analítica de números %%%%%%%%%%%%%%%%%%%%%%%%%%%%%%
%%%%%%%%%%%%%%%%%%%%%%%%%%%%%%%%%%%%%%%%%%%%%%%%%%%%%%%%%%%%%%
\newpage
\subsection*{MA-0609 Teoría analítica de números}
\addcontentsline{toc}{subsection}{MA-0609 Teoría analítica de números}

\hypertarget{MA0609}{}
\begin{refsection}
\begin{table}[H]
\centering{
\begin{tabular}{|ll|}
\hline
%\multicolumn{2}{|c|}{\textbf{Nivel de virtualidad del curso:} Presencial}\\ \hline
\textbf{Año:} IV  & \textbf{Requisitos:} MA-0496 Teoría de Números \\
& y MA-0702 Análisis Complejo \\
\textbf{Ciclo:} VIII  & \textbf{Correquisitos:} Ninguno \\
\textbf{Tipo de curso:} Teórico & \textbf{Horas presenciales por semana:} 5 \\
\textbf{Créditos:} 4 & \textbf{Horas de trabajo independiente por semana:} 7 \\
\textbf{Virtualidad:} P  &  \\ \hline
\end{tabular}
}
\end{table}

\subsubsection*{Descripción del curso}

Este curso introduce a la persona estudiante a las ideas y métodos fundamentales de la Teoría Analítica de Números. El objetivo central es estudiar problemas aritméticos mediante herramientas provenientes del análisis real y complejo, con especial énfasis en el comportamiento promedio de funciones aritméticas y en la distribución de los números primos.

El curso comienza con el estudio de funciones aritméticas, convolución de Dirichlet, funciones sumatorias y órdenes promedio, junto con técnicas básicas para la estimación de sumas. Posteriormente se introducen las series de Dirichlet y las herramientas analíticas que permiten recuperar información sobre sus coeficientes, incluyendo la suma parcial de Abel, la transformada de Mellin y la fórmula de Perron.

Una parte central del curso se dedica al estudio de la función zeta de Riemann y de las funciones $L$ de Dirichlet, incluyendo sus propiedades de convergencia, productos de Euler, continuación analítica, ecuaciones funcionales y distribución de sus ceros. Estas herramientas se aplican al estudio de la distribución de los números primos, incluyendo el Teorema de los Números Primos y los primos en progresiones aritméticas.

El curso también busca presentar a la persona estudiante algunos de los grandes problemas abiertos de la teoría, especialmente la Hipótesis de Riemann y sus generalizaciones, así como resultados clásicos que muestran la estrecha relación entre la localización de los ceros de funciones zeta y funciones $L$ y la distribución de los números primos.

\subsubsection*{Objetivos}

\textbf{General:} Introducir a la persona estudiante a las principales ideas, técnicas y resultados de la Teoría Analítica de Números, enfatizando el uso de herramientas del análisis real y complejo para estudiar funciones aritméticas y la distribución de los números primos. \

\textbf{Específicos:}

\begin{enumerate}
\item Estudiar las propiedades fundamentales de las funciones aritméticas y sus funciones sumatorias.
\item Utilizar técnicas de sumación y estimación para determinar órdenes promedio y comportamientos asintóticos.
\item Comprender la relación entre funciones aritméticas, series de Dirichlet y productos de Euler.
\item Aplicar herramientas del análisis complejo, incluyendo integración de contorno y transformadas integrales, al estudio de problemas aritméticos.
\item Estudiar las propiedades analíticas fundamentales de la función zeta de Riemann y de las funciones $L$ de Dirichlet.
\item Comprender el papel de los ceros de funciones zeta y funciones $L$ en problemas de distribución de números primos.
\item Estudiar el Teorema de los Números Primos y el teorema de Dirichlet sobre primos en progresiones aritméticas.
\item Conocer algunos de los principales problemas abiertos y conjeturas de la Teoría Analítica de Números.
\item Comunicar argumentos y resultados matemáticos de manera clara y rigurosa en forma oral y escrita.
\end{enumerate}

\subsubsection*{Contenidos}

\begin{enumerate}

\item \textbf{Funciones aritméticas y convolución de Dirichlet (2 semanas):}
\begin{enumerate}
\item Funciones aritméticas multiplicativas y completamente multiplicativas.
\item Convolución de Dirichlet y estructura algebraica de las funciones aritméticas.
\item Funciones clásicas: función de M\"obius, función de Euler, funciones divisor y funciones de von Mangoldt.
\item Fórmulas de inversión de M\"obius y aplicaciones.
\end{enumerate}

\item \textbf{Funciones sumatorias y órdenes promedio (2 semanas):}
\begin{enumerate}
\item Funciones sumatorias y órdenes promedio de funciones aritméticas.
\item Notaciones asintóticas y técnicas elementales de estimación.
\item Suma parcial de Abel y sus aplicaciones.
\item Método de la hipérbola de Dirichlet y otras técnicas elementales de sumación.
\item Ejemplos de estimaciones asintóticas para funciones aritméticas clásicas.
\end{enumerate}

\item \textbf{Series de Dirichlet y herramientas de análisis complejo (2 semanas):}
\begin{enumerate}
\item Series de Dirichlet: convergencia, convergencia absoluta y propiedades analíticas básicas.
\item Relación entre convolución de Dirichlet y productos de series de Dirichlet.
\item Productos de Euler.
\item Transformada de Mellin y fórmula de inversión.
\item Fórmula de Perron y recuperación de funciones sumatorias.
\item Aplicaciones de la integración de contorno a estimaciones aritméticas.
\end{enumerate}

\item \textbf{La función zeta de Riemann (2 semanas):}
\begin{enumerate}
\item Definición mediante serie de Dirichlet y producto de Euler.
\item Continuación meromorfa y comportamiento en el punto singular.
\item Ecuación funcional y función zeta completada.
\item Ceros triviales y no triviales.
\item Producto de Hadamard y relación entre ceros y comportamiento analítico.
\end{enumerate}

\item \textbf{Ceros de la función zeta y distribución de los números primos (2 semanas):}
\begin{enumerate}
\item Funciones de Chebyshev y formulaciones equivalentes del Teorema de los Números Primos.
\item Regiones libres de ceros y no anulación sobre la recta de parte real uno.
\item Teorema de los Números Primos.
\item Fórmulas explícitas que relacionan números primos y ceros de la función zeta.
\item Conteo de ceros y fórmula de Riemann--von Mangoldt.
\item Hipótesis de Riemann y algunas de sus consecuencias.
\item Enunciado y discusión de resultados más refinados sobre la distribución de los ceros.
\end{enumerate}

\item \textbf{Caracteres y funciones $L$ de Dirichlet (2 semanas):}
\begin{enumerate}
\item Caracteres de Dirichlet y relaciones de ortogonalidad.
\item Caracteres primitivos e inducidos.
\item Sumas de Gauss.
\item Funciones $L$ de Dirichlet y sus productos de Euler.
\item Continuación analítica y ecuaciones funcionales.
\item Ceros de funciones $L$ y generalización de la Hipótesis de Riemann.
\end{enumerate}

\item \textbf{Primos en progresiones aritméticas (2 semanas):}
\begin{enumerate}
\item No anulación de funciones $L$ en el punto uno.
\item Teorema de Dirichlet sobre primos en progresiones aritméticas.
\item Distribución asintótica de primos en progresiones aritméticas.
\item Influencia de los ceros de las funciones $L$ sobre los términos de error.
\item Enunciados de algunos resultados clásicos sobre regiones libres de ceros y ceros excepcionales.
\end{enumerate}

\item \textbf{Problemas abiertos y desarrollos posteriores (1 semana):}
\begin{enumerate}
\item La Hipótesis de Riemann y la Hipótesis de Riemann Generalizada.
\item Problemas abiertos relacionados con la distribución de los números primos.
\item Resultados de densidad de ceros y consecuencias para la distribución de primos, a nivel de enunciados.
\item Introducción a otras direcciones de la Teoría Analítica de Números, tales como métodos de criba, sumas exponenciales y funciones $L$ más generales.
\end{enumerate}

\end{enumerate}

\subsubsection*{Metodología}

\noindent \textbf{Lineamientos metodológicos:} La persona docente a cargo de este curso debe respetar las siguientes pautas:
\begin{enumerate}
\item El curso debe combinar el estudio riguroso de resultados fundamentales con el desarrollo de técnicas concretas para estimar sumas y funciones aritméticas.
\item Las herramientas de análisis complejo deben introducirse en conexión directa con problemas aritméticos, aprovechando los conocimientos adquiridos previamente en el curso de Variable Compleja.
\item Se debe enfatizar la relación entre el comportamiento analítico de series de Dirichlet y funciones $L$ y la información aritmética contenida en sus coeficientes.
\item El estudio de la función zeta de Riemann y las funciones $L$ de Dirichlet debe incluir tanto resultados demostrados en el curso como una discusión de resultados más profundos y problemas abiertos.
\item Se debe distinguir claramente entre los resultados que se demuestran en el curso y aquellos que se presentan únicamente a nivel de enunciado o discusión.
\end{enumerate}

\textbf{Sugerencias metodológicas:} Adicionalmente, se tienen las siguientes recomendaciones para la persona docente a cargo de este curso:
\begin{enumerate}
\item Utilizar ejemplos concretos de funciones aritméticas para motivar las técnicas generales antes de formular resultados abstractos.
\item Desarrollar ejemplos detallados de suma parcial, método de la hipérbola, series de Dirichlet, fórmula de Perron y desplazamiento de contornos.
\item Complementar algunos temas con experimentación computacional, por ejemplo para comparar funciones contadoras de primos con sus aproximaciones asintóticas o para visualizar ceros de la función zeta.
\item Relacionar los principales resultados con su desarrollo histórico, especialmente los trabajos de Euler, Dirichlet, Riemann, Hadamard, de la Vallée Poussin y otros autores fundamentales.
\item En la evaluación, se sugiere combinar problemas de demostración con ejercicios de estimación y aplicación de técnicas analíticas.
\item Promover la lectura de secciones seleccionadas de textos más avanzados con el objetivo de familiarizar a las personas estudiantes con resultados que conducen hacia la investigación contemporánea.
\item Promover la comunicación clara y rigurosa de argumentos e ideas matemáticas de manera oral y escrita.
\end{enumerate}

\subsubsection*{Guía para las referencias}

El libro de Tom M. Apostol \cite{Apo76} es una referencia principal particularmente adecuada para la primera parte del curso. Su tratamiento de funciones aritméticas, convolución de Dirichlet, órdenes promedio, series de Dirichlet y productos de Euler proporciona una introducción accesible y rigurosa a las técnicas fundamentales de la Teoría Analítica de Números. También contiene un tratamiento introductorio de la función zeta de Riemann, las funciones $L$ de Dirichlet y el Teorema de los Números Primos.

El libro de Marius Overholt \cite{Ove15} constituye otra posible referencia principal para el curso y permite avanzar de manera natural hacia técnicas analíticas más sofisticadas, incluyendo la fórmula de Perron, integración de contorno, funciones zeta y funciones $L$. El libro de M. Ram Murty \cite{Mur08}, organizado alrededor de problemas y sus soluciones, es especialmente útil como fuente de ejercicios y para complementar los temas estudiados en clase.

El libro clásico de Harold Davenport \cite{Dav00}, en su tercera edición revisada por Hugh L. Montgomery, es una referencia fundamental para el estudio de funciones $L$ de Dirichlet y la distribución de primos en progresiones aritméticas. Finalmente, el libro de Hugh L. Montgomery y Robert C. Vaughan \cite{MV07} ofrece un tratamiento más avanzado y sistemático de la teoría multiplicativa de números, incluyendo series de Dirichlet, el Teorema de los Números Primos, propiedades analíticas de funciones zeta y funciones $L$, fórmulas explícitas y distribución de sus ceros. Se recomienda particularmente como referencia para ampliar los temas tratados y como puente hacia cursos de posgrado e investigación.

\nocite{Apo76}
\nocite{Dav00}
\nocite{Ove15}
\nocite{Mur08}
\nocite{MV07}

\printcursosbibliography
\end{refsection}
